\SourceLeaf{90}%

Als Endresultat haben wir also

\[\begin{aligned}
fx={}&ft+\psi t\mathord{.}\frac{dft}{dt}\mathord{.}u
+\frac{1}{2}\frac{d(\psi t)^2\mathord{.}\frac{dft}{dt}}{dt}\mathord{.}uu\\
&+\frac{1}{2.3}\frac{dd(\psi t)^3\mathord{.}\frac{dft}{dt}}{dt^2}\mathord{.}u^3
+\frac{1}{2.3.4}\frac{d^3(\psi t)^4\frac{dft}{dt}}{dt^3}u^4+\text{u.\,s.\,w.},
\end{aligned}\]

\noindent welches der \textsc{Lagrang}ische Lehrsatz in seiner ganzen Allgemeinheit ist.

1847 Mai 13.
